\section{An exact reduction for massless banana integrals}
\label{app:banana}
\subsection{Definition and convolution}
Consider the massless two-point family
\begin{equation}
 B_L(\boldsymbol a;p^2,d)=
 \int\prod_{j=1}^{L}\frac{d^d\ell_j}{i\pi^{d/2}}
 \frac{1}{\prod_{j=1}^{L}(\ell_j^2+i0)^{a_j}
 [(p-\sum_{j=1}^{L}\ell_j)^2+i0]^{a_{L+1}}}.
 \label{eq:banana-family}
\end{equation}
Here $L\geq1$, all $a_i$ are positive integers, and $p^2\ne0$. The family
has two vertices joined by $L+1$ massless lines and is often called a banana
or watermelon diagram. The implemented specialization covers the displayed
scalar integrals. Auxiliary denominator entries may be present in the Kira
definition, but their target indices must be zero. The proof does not cover
arbitrary numerator insertions or other propagator structures.

The one-loop convolution underlying this calculation is standard~\cite{Grozin2004Banana}.
For completeness, write its Euclidean form as
\begin{equation}
 E_1(a,b;p_E^2)=\int\frac{d^d k_E}{\pi^{d/2}}
                 (k_E^2)^{-a}[(p_E-k_E)^2]^{-b}.
\end{equation}
Here $k_E$ and $p_E$ are Euclidean loop and external momenta, and $a,b$
are the two denominator powers. The Euler Gamma function is
$\Gamma(z)=\int_0^\infty \eta^{z-1}e^{-\eta}d\eta$ for $\operatorname{Re}z>0$.
Changing the integration variable to $\eta=tD$ gives the Schwinger
representation $D^{-a}=\Gamma(a)^{-1}\int_0^\infty t^{a-1}e^{-tD}dt$
for positive $D$ in its convergence domain. Applying this representation
to both denominators and completing the square gives
\begin{equation}
 t k_E^2+u(p_E-k_E)^2
 =(t+u)\left(k_E-\frac{u}{t+u}p_E\right)^2
  +\frac{tu}{t+u}p_E^2.
\end{equation}
For $\alpha>0$, squaring $\int_{\mathbb R}e^{-\alpha x^2}dx$ gives
a two-dimensional Gaussian integral. Polar coordinates evaluate it as
$2\pi\int_0^\infty r e^{-\alpha r^2}dr=\pi/\alpha$.
The positive square root and the product over momentum components give
the normalized Gaussian momentum integral $(t+u)^{-d/2}$, continued in
$d$ from integer dimension. Thus the Schwinger parameters $t,u>0$ give
\begin{equation}
 E_1=\frac{1}{\Gamma(a)\Gamma(b)}
 \int_0^\infty dt\,du\,
 t^{a-1}u^{b-1}(t+u)^{-d/2}
 \exp\left[-\frac{tu}{t+u}p_E^2\right].
\end{equation}
Set $t=\lambda x$, $u=\lambda(1-x)$, with Jacobian $\lambda$,
$\lambda>0$, and $0<x<1$. The $\lambda$ integral gives
$\Gamma(a+b-d/2)[x(1-x)p_E^2]^{d/2-a-b}$. The remaining integral is
$\int_0^1 x^{d/2-b-1}(1-x)^{d/2-a-1}dx$. To evaluate it, apply the same
sum-and-fraction change of variables to the two integrals defining
$\Gamma(z)\Gamma(w)$. The sum integral gives $\Gamma(z+w)$, which yields
\begin{equation}
 \int_0^1 x^{z-1}(1-x)^{w-1}dx
 =\frac{\Gamma(z)\Gamma(w)}{\Gamma(z+w)}
 \quad (\operatorname{Re}z,\operatorname{Re}w>0).
\end{equation}
Taking $z=d/2-b$ and $w=d/2-a$ then gives
\begin{equation}
 E_1=(p_E^2)^{d/2-a-b}G(a,b),\qquad
 G(a,b)=\frac{\Gamma(a+b-d/2)\Gamma(d/2-a)\Gamma(d/2-b)}
 {\Gamma(a)\Gamma(b)\Gamma(d-a-b)}.
 \label{eq:banana-convolution}
\end{equation}
The parameter integrations first hold in a common convergence region.
Analytic continuation defines the dimensionally regulated expression
outside that region, as in the usual one-loop formula~\cite{Grozin2004Banana}.

Let $E_L$ be the Euclidean version of Eq.~\eqref{eq:banana-family}, with
measure $\prod_j d^d\ell_{j,E}/\pi^{d/2}$, and write $n=L+1$ and
$A=\sum_{i=1}^{n}a_i$. Integrating the last loop combines
the last two indices into $b=a_L+a_{L+1}-d/2$. Applying
Eq.~\eqref{eq:banana-convolution} and then the $(L-1)$-loop formula cancels
the factors $\Gamma(b)$ and $\Gamma(d-a_L-a_{L+1})$ between the two stages.
Induction from $L=1$ therefore gives
\begin{equation}
 E_L(\boldsymbol a;p_E^2,d)=
 (p_E^2)^{Ld/2-A}
 \frac{\Gamma(A-Ld/2)}{\Gamma(nd/2-A)}
 \prod_{i=1}^{n}\frac{\Gamma(d/2-a_i)}{\Gamma(a_i)}.
 \label{eq:banana-gamma}
\end{equation}
Wick rotation to the denominator convention in Eq.~\eqref{eq:banana-family}
gives one factor of $i$ per loop measure, cancelling its normalization,
and changes the sign of each quadratic denominator. It therefore
replaces the power factor by $(-1)^A(-p^2-i0)^{Ld/2-A}$. The remaining
Gamma factors stay unchanged. This sign convention matters when reducing
an integral with an additional denominator power.

\subsection{Rational coefficients in a single master}
Choose $M_L=B_L(1,\ldots,1;p^2,d)$ and define $b_i=a_i-1\geq0$ and
$K=\sum_i b_i=A-n$. The Gamma function recurrence gives
\begin{equation}
 \frac{\Gamma(z+K)}{\Gamma(z)}=(z)_K,\qquad
 \frac{\Gamma(z)}{\Gamma(z-K)}=(-1)^K(1-z)_K,
\end{equation}
where $(z)_K=\prod_{j=0}^{K-1}(z+j)$ and $(z)_0=1$. Integration by parts
in the defining Gamma integral gives $\Gamma(z+1)=z\Gamma(z)$ when the
boundary terms vanish; analytic continuation extends the relation.
Applying this recurrence repeatedly gives the first identity. Reversing the
finite product and factoring out $-1$ from each factor gives the second.
In the ratio of Eq.~\eqref{eq:banana-gamma} to its corner value, the
dimension-independent index factors give $\Gamma(1)/\Gamma(a_i)=1/b_i!$.
The shifted factors satisfy
\begin{equation}
 \prod_i\frac{\Gamma(d/2-1-b_i)}{\Gamma(d/2-1)}
   =\frac{(-1)^K}{\prod_i(2-d/2)_{b_i}}.
\end{equation}
This sign cancels the sign from the inverse Gamma ratio. The remaining
Minkowski power factor is
$(-1)^K(-p^2-i0)^{-K}=(p^2+i0)^{-K}$ because $K$ is an integer.
Thus the required reduction is
\begin{equation}
 \boxed{\displaystyle
 B_L(\boldsymbol a;p^2,d)=
 \frac{(n-Ld/2)_K\,(n+1-nd/2)_K}
 {(p^2)^K\prod_{i=1}^{n}(2-d/2)_{b_i}\,b_i!}\,M_L.}
 \label{eq:banana-ratio}
\end{equation}
We understand the rational coefficient by meromorphic continuation in $d$
and generic nonzero $p^2$. The corner value is generically nonzero, and
Eq.~\eqref{eq:banana-ratio} expresses every scalar target in this class
through it. This establishes a one-element basis for the stated scalar
class without a new master-discovery calculation.

The implementation forms finite products of linear factors in $d$ and
cancels common factors with exact rational arithmetic before formatting
the result. It uses no trained model to predict a coefficient. The domain
check verifies the physical denominator momenta, zero masses, the single
external squared momentum, positive physical indices, and zero auxiliary
indices. If the caller supplies a basis, that basis must contain $M_L$.
The implementation caps the loop count at 64 and the total extra powers
at 256 to bound work; the mathematical identity holds for every positive
integer loop count under the stated assumptions.

\subsection{Independent comparisons}
The native comparison uses three targets at each of one through six loops,
with rational squared momentum $p^2=-1651$, and repeats the calculation
three times. A second comparison keeps $p^2$ symbolic at one through four
loops. The requested indices raise one physical power by one or two.
Kira receives the corner basis, generates native IBP equations with all ML
disabled, and returns rational coefficients. Exact subtraction checks these
coefficients against Eq.~\eqref{eq:banana-ratio}. Both reported timings include
their coefficient check; the direct formula also includes its domain check.
Shared Python imports precede these timings.

Eight initial higher-loop requests fail, including corner-basis rejections
and an allocation failure in the six-loop preparation. We retain them and
make a separate fixed-seed comparison. Let $r$ be the
largest sum of positive target indices, $s$ the largest numerator rank,
and $q$ the largest number of extra denominator powers. Development at
$p^2=-1649$ compares seed bounds $(r-1,s,q-1)$, $(r,s,q)$,
$(r,s+1,q)$, and $(r+1,s+1,q+1)$; the positive-power limit never falls
below the number of physical denominators. Symbolic development changes
the raised propagator positions at the same index profile. For each loop
count and kinematic form, the fastest successful complete development
calculation fixes the seed choice. A new comparison at $p^2=-1655$, or
symbolic $p^2$ with the new index profile, uses three repetitions of that
fixed choice. Timed failures remain separate from matched speedup ratios.

Further checks use one through twelve loops and loop counts 16, 24, and 32,
with eight distributions of one through eight extra powers at each loop
count. The direct rational result is compared with the uncancelled Gamma
ratio from Eq.~\eqref{eq:banana-gamma} at $d=257/101$ and $d=401/97$.
The evaluation uses 70-digit arithmetic and a relative error threshold of
$10^{-55}$ with denominator $\max(1,|r|)$, where $r$ is the Gamma-ratio
value. These numerical checks supplement the derivation. Native timing
ratios are reported only where the native calculation completes; no native
speedup is assigned to the higher-loop cases checked by the formula alone.

The combined-interface application check uses a fresh process for each
request. Four scalar banana requests at one, four, eight, and twelve loops
use an independent symbolic implementation of Eq.~\eqref{eq:banana-ratio}
as their reference. It constructs rising factorials directly and does not
call the live factor-cancellation routine. The native application requests
use the independently checked IBP references of the target-batch study.
All reference construction precedes application timing. Both comparison
strategies use the same analytic route; their difference on native requests
comes from the frozen learned or fixed reduction choices.
Exact rational checks at $d=257/101$ also test the implemented power limit:
one- and 64-loop targets have 256 extra powers on one denominator. A direct
product calculation checks the factored result without expanding its
polynomials. Further checks verify the corner identity and reject requests
beyond the implemented loop or power limits.
